\subsection{Grundy y Sprague-Grundy}\label{rec:grundy}

Un juego imparcial se parte en componentes independientes. El SG de una
posición es el \codigo{mex} de los SG de sus movimientos. Si el XOR de los SG
de las componentes es distinto de 0, gana el primero.

\tablafunc{
\codigo{mex(v)} & menor entero no negativo que no está en \codigo{v} & $O(k)$ \\
\codigo{subgrundy} & SG de una pila de 0 a \codigo{n} & $O(nm)$ \\
\codigo{Grundy::calc()} & SG de un DAG (o \texttt{\{\}} si hay ciclo) & $O(V+E)$ \\
}

\copiar{el bloque \texttt{COPIA GRUNDY}.}
\lstinputlisting{codigo/grundy.cpp}

\ojo{Para varias pilas o componentes, combina con XOR:
\texttt{g = g1 \textasciicircum{} g2 \textasciicircum{} ...}. Si
\codigo{g != 0} gana el primero. Un estado terminal tiene SG 0.}